\\maketitle
\\begin{abstract}
    本文针对2022年CUMCM B题新能源汽车充电策略优化问题，建立了基于混合整数规划的充电调度模型。首先分析用户充电行为特征和电网负荷特性；其次构建以最小化充电成本、最大化用户满意度、均衡电网负荷为目标的多目标优化模型；最后采用NSGA-II算法求解。
    
    \\textbf{针对问题一}，我们分析了用户充电行为。通过问卷调查和数据分析，建立充电需求预测模型，预测结果显示早晚高峰充电需求占总需求的65\\%。
    
    \\textbf{针对问题二}，我们建立了单目标充电调度模型。决策变量为$x_{i,t}$表示车辆$i$在时刻$t$是否充电，目标函数：
    $$\\min Z = \\sum_{i,t} (C_{elec} \\cdot P_{i,t} \\cdot x_{i,t} + C_{delay} \\cdot d_{i,t})$$
    
    \\text\"针对问题三}，我们扩展为多目标优化模型。考虑电网峰谷电价、电池损耗成本、用户等待时间等多因素，采用NSGA-II算法求解得到Pareto最优解集。
    
    本研究为电动汽车智能充电管理提供了科学决策支持。
    
    \\keywords{电动汽车\quad 充电调度\quad 混合整数规划\quad NSGA-II\quad 智能电网}
\\end{abstract}
